\par\bigskip
\SourceLeaf{489}%
\pdfbookmark[0]{Entwurf zur Gradmessung. 1819.}{entwurf-gradmessung-1819}
\begin{center}\large [ENTWURF ZUR GRADMESSUNG.]\end{center}
\par
— — —

Durch ein Schreiben des Professors \textsc{Schumacher} bin ich benachrichtigt, dass der Dänische Gesandte in London Graf \textsc{Bourke}, bei Gelegenheit von des erstern Aufenthalt daselbst, mit dem Grafen \textsc{Münster} über den Nutzen der Fortsetzung der Dänischen Gradmessung durch das Königreich Hannover gesprochen, und dass letzterer verlangt habe, dass ich über diesen Gegenstand an ihn schreiben möchte. Wenn ich diese Anzeige wie einen Befehl habe betrachten müssen, so verpflichtet mich das warme Interesse, welches Ew. — — an dieser wissenschaftlichen Angelegenheit genommen, und bereits im vorigen Jahre, durch die die Möglichkeit einer solchen Fortsetzung sichernden Maassregeln bethätigt haben, Ew. — — sofort von diesem Umstande zu benachrichtigen, und bürgt mir für die gütige Aufnahme meiner Bitte, die Sache nach Möglichkeit in London zu unterstützen. Ich darf mich um so weniger scheuen, diese Bitte zu thun, da ich dabei bloss aus reiner Liebe zur Wissenschaft handle, und ohne diese und ohne die lebhafte Überzeugung von der hohen Wichtigkeit einer solchen Operation, nach meiner individuellen Neigung, eine längere Entfernung von den rein astronomischen Arbeiten eher als ein Opfer betrachten würde. Auch der ehrwürdige Sir \textsc{Joseph Banks} hat, wie mir Professor \textsc{Schumacher} schreibt, die Idee der weitern Fortsetzung dieser Gradmessung auf dem Continent, mit grosser Wärme ergriffen, und alle Mitwirkung versprochen.

Der Professor \textsc{Schumacher} wird bei seiner Durchreise durch Hannover Ew. — — selbst gesagt haben, dass der Zweck seiner Reise nach England die Empfangnahme des berühmten grossen Zenithsectors war, welcher bei der englischen Gradmessung gebraucht ist, und ihm für seine astronomischen Beobachtungen geliehen wird. Er wird sogleich nach seiner Zurückkunft, diese Beobachtungen am südlichen Endpunkte seiner Messungen in Lauenburg anfangen. Bei der Aussicht zu einer Fortsetzung der Dänischen Gradmessung durch das Königreich Hannover würde es ohne Zweifel in vielfacher Hinsicht für beide Gradmessungen sehr wünschenswerth sein, wenn ich \textit{diese} Beobachtungen mit ihm gemeinschaftlich machte. \textit{Ohne} diese Vereinigung ist der Fall leicht denkbar, dass in Zukunft diese oder jene Discordanzen oder Zweifel eintreten könnten, die sich dann nur schwer würden aufklären lassen, und denen \textit{durch} dieselbe vorgebeugt werden würde. So sind ja gerade über die astronomischen Beobachtungen, welche am nordlichen Endpunkte der französischen Gradmessung in Dünkirchen angestellt waren, späterhin Bedenklichkeiten entstanden, derentwegen es nöthig gefunden wurde, dass im vorigen Herbst mit dem gedachten Zenithsector neue Beobachtungen durch englische und französische Commissarien gemeinschaftlich, angestellt wurden; und wie ich höre, soll eben dieser Zenithsector, sobald ihn Professor \textsc{Schumacher} wieder abgegeben haben wird, zu einer nochmaligen Wiederholung der Beobachtungen auf den Balearischen Inseln, dahin eingeschifft werden. Nur durch eignes Beobachten an diesem Instrumente, und durch Vergleichung mit den Erfahrungen, die ich an dem \textsc{Repsold}schen Meridiankreise gemacht habe, und noch mehr mit denen, welche ich mit dem \textsc{Reichenbach}schen Kreise in Zukunft machen werde, würde ich im Stande sein, zu beurtheilen, ob eine künftige Zuziehung eben jenes Zenithsectors für die Beobachtungen im Königreich Hannover nöthig oder vorzüglich wünschenswerth sein wird. Und selbst den Fall angenommen, dass die Fortsetzung der Gradmessung nicht zur Ausführung käme, würden die von mir an dem Zenithsector, einem in seiner Art einzigen Instrumente, gemachten Erfahrungen für meine künftigen Beobachtungen an den Meridian-Instrumenten der hiesigen Sternwarte nicht ohne wesentlichen Nutzen sein. Als einen Beweis, wie wichtig die Beobachtungen in Lauenburg geachtet werden, darf ich noch anführen, dass der Doctor \textsc{Olbers} mir von seiner Absicht geschrieben hat, nach Lauenburg kommen zu wollen, um das Instrument kennen zu lernen und den Beobachtungen beizuwohnen.
